\section{Post-training operating system 与开放资产栈}

访谈最后把竞争重点从一次性 pre-training 转向持续 post-training。Lample 的判断是，预训练架构在较长时间内保持相似，而 post-training 仍有大量可创新空间。这个判断不是说预训练不重要，而是说产品差异越来越依赖 environment、verifier、数据生成、tool use、serving 和快速实验 pipeline 的联合质量。

\subsection{为什么 post-training 不是一个算法名}

\term{post-training} 泛指 base model 之后的行为适配，包括 supervised fine-tuning、preference optimization、RL、distillation、tool-use training 和持续 evaluation。任何一个环节若没有版本、指标和回滚，系统都难以快速实验。因此可把 post-training 看成一个 operating system：上层接收产品任务，下层连接训练与 serving，横向提供数据、评测和治理。

\lecturefigure{14-post-training-os.png}{Post-training 的竞争力来自整套可重复迭代系统。}{本地字幕 37:10--39:25；概念重绘。}

\begin{knowledgebox}{读图：灵活性为什么可能胜过单纯规模}
Online 或高频 post-training 需要短实验周期、稳定环境、自动评测和清晰 ownership。超大组织若 pipeline 分散、审批和数据接口迟缓，额外 compute 不能自动转化为迭代速度；小团队若缺少纪律，同样会被不可复现实验拖垮。优势来自“足够计算 + 高信息反馈 + 低摩擦执行”。
\end{knowledgebox}

\subsection{开放模型如何进入商业栈}

当其他实验室发布强 open-weight 模型，解决方案公司可以把它当作已经支付过 pre-training compute 的资产，再进行领域定制、工具集成和私有部署。价值并没有消失，而是从 base model 稀缺性向数据、workflow、可靠性和控制上移。

\lecturefigure{16-open-asset-stack.png}{开放权重是产品栈的输入，不是完整价值交付。}{本地字幕 39:30--42:30；Mistral 7B 官方开放许可与部署资料。}

\begin{importantbox}{开放资产栈的工程检查表}
\begin{enumerate}
  \item 验证 license、模型卡、数据和已知限制；
  \item 在目标硬件和 workload 上建立 baseline；
  \item 建立领域数据、评测与安全边界；
  \item 将工具、检索、权限和业务状态纳入端到端测试；
  \item 为 base model 更新准备兼容性测试和回滚路径。
\end{enumerate}
\end{importantbox}

\begin{warningbox}{不要把“可复用”误写成“没有护城河”}
开放模型降低了基础能力获取成本，也提高了客户对迁移和控制的预期。真正的竞争壁垒可能转向高质量反馈、领域数据权利、workflow integration、评测资产、分发和运营可靠性；这些仍然需要长期投入。
\end{warningbox}

\subsection{本章小结}

Post-training 是连接产品任务与模型行为的操作系统。开放资产让团队复用公共 compute，但差异化必须通过数据、工具、部署和反馈闭环重新建立。

